\documentclass[11pt]{article}
\usepackage[utf8]{inputenc}
\usepackage{graphicx}
\usepackage[margin=0.5in, includeheadfoot, footskip=133pt]{geometry}
\usepackage{xcolor}
\usepackage{hyperref}
\usepackage{fancyhdr}
\usepackage{enumitem}
\usepackage{titlesec}
\usepackage{amsmath}
\usepackage{amssymb}
\usepackage{tikz}
\usepackage{unicode-math}
\usetikzlibrary{arrows,arrows.meta,calc,positioning}

% Define Caltech orange color
\definecolor{caltechorange}{RGB}{255, 108, 12}

% Define tightlist command
\providecommand{\tightlist}{%
  \setlength{\itemsep}{0pt}\setlength{\parskip}{0pt}}

% Section formatting
\titleformat{\section}
  {\normalfont\Huge\bfseries}{}{0em}{}
\titlespacing*{\section}
  {0pt}{3.5ex plus 1ex minus .2ex}{2.3ex plus .2ex}

% Subsection formatting - no numbers
\titleformat{\subsection}
  {\normalfont\Large\bfseries}{}{0em}{}
\titlespacing*{\subsection}
  {0pt}{3.25ex plus 1ex minus .2ex}{1.5ex plus .2ex}

% Subsubsection formatting - no numbers  
\titleformat{\subsubsection}
  {\normalfont\large\bfseries}{}{0em}{}
\titlespacing*{\subsubsection}
  {0pt}{3.25ex plus 1ex minus .2ex}{1.5ex plus .2ex}

% Page style
\pagestyle{fancy}
\fancyhf{}

% Header setup
\fancyhead[L]{%
    \includegraphics[height=1cm]{templates/caltech_logo.png}
}
\fancyhead[C]{%
    {\color{gray}\small Center for Technology \\ Management Education}
}
\fancyhead[R]{%
    {\color{gray}\small\href{https://ctme.caltech.edu}{ctme.caltech.edu}}
}

% Add padding after header
\setlength{\headheight}{50pt}
\setlength{\headsep}{2em}

% Force fancy style on first page
\AtBeginDocument{\thispagestyle{fancy}}

% Footer

\cfoot{
\begin{tabular}{p{8in}}
  1200 E California Blvd • MC 121-79 • Pasadena, California 91125 • 626.395.4042 • execed@caltech.edu
  \\ \ 
  \\ \ 
\end{tabular}
}


% List formatting
\setlist{nosep}  % Remove vertical spacing between list items
\setlist[itemize]{leftmargin=*}  % Align bullet points with left margin

% Hyperref setup
\hypersetup{
    colorlinks=true,
    linkcolor=caltechorange,
    filecolor=caltechorange,
    urlcolor=caltechorange,
    citecolor=caltechorange
}

% Document
\begin{document}

\section{Linear Algebra and Estimation
Theory}\label{linear-algebra-and-estimation-theory}

\paragraph{Course Details:}

\textbf{Title:} Linear Algebra and Estimation Theory\\
\textbf{Duration:} 10 weeks\\
\textbf{Lecturers:} TBD\\
\textbf{Format:} Live-remote, twice weekly 90-minute lectures\\
\textbf{Start:} January 7, 2025

\paragraph{Course Description:}

This course covers. foundational concepts in linear algebra as they
relate to estimation techniques. It equips participants with the
mathematical tools necessary for applications in signal processing,
communications, and data analysis. Participants will delve into
essential topics such as random vectors, covariance matrices, linear
transformations, and Gaussian random variables. Through interactive
exercises and real-world examples, students will gain practical insights
into the analytical frameworks critical to engineering and applied
sciences.

\begin{itemize}
\tightlist
\item
  \textbf{Note:} I am assuming that you are using an LLM to help you
  with your work in this course.
\end{itemize}

\paragraph{Learning Outcomes:}

\begin{itemize}
\tightlist
\item
  Master the principles of vector spaces, matrices, and linear
  transformations
\item
  Understand and apply the properties of random vectors and covariance
  matrices
\item
  Derive and utilize correlation and cross-correlation matrices in data
  analysis
\item
  Perform linear transformations and interpret their effects on data
\item
  Apply the Central Limit Theorem and Law of Large Numbers in estimation
  contexts
\item
  Calculate and interpret expectations, variances, and covariance for
  multivariable functions
\item
  Design and evaluate linear estimators using minimum mean square error
  (MMSE) techniques
\item
  Develop analytical skills for practical applications in signal
  processing and communications
\end{itemize}

\end{document}
